% =====================================================================
%  Response to Reviewers — template (NeurIPS / E&D style)
%  Compile: pdflatex response_template.tex
%  Convention: quote the reviewer comment, then respond; cite line numbers.
% =====================================================================
\documentclass[11pt]{article}
\usepackage[utf8]{inputenc}
\usepackage[T1]{fontenc}
\usepackage[margin=1in]{geometry}
\usepackage{xcolor}
\usepackage{parskip}
\usepackage{enumitem}

\newcommand{\rev}[1]{\textbf{\color{blue}Reviewer:} \emph{``#1''}}
\newcommand{\resp}{\textbf{\color{red}Response:}}
\newcommand{\chg}[1]{\textbf{Change:} #1}

\begin{document}

\begin{center}
{\Large\bfseries Response to Reviewers}\\[2pt]
{\large Queyi: Making Verification Capability Independently Auditable}\\[2pt]
\textbf{Paper \#XXXX} \quad (Anonymous submission)
\end{center}

\noindent
We thank all reviewers for their careful reading. We address every point below; reviewer comments
are in \emph{italic}, our responses follow. Section/line numbers refer to the revised manuscript.

\section*{Summary of changes}
\begin{itemize}[leftmargin=1.4em]
  \item Added the budget-matched baseline results (B1/B2/B3) and the A--F ablation in \S6.
  \item Expanded the holdout/corpus sample sizes (holdout $\ge$30 measurable; corpus $\ge$60).
  \item Added Clopper--Pearson CIs to all headline numbers.
  \item Clarified the environment dependency (WSL) and the \texttt{UNVERIFIED} protocol.
\end{itemize}

\section*{Reviewer 1}

\rev{The paper claims a new evaluation paradigm, but the empirical support is thin.}
\resp We agree that the initial submission was thin on baselines. We have now run the three
budget-matched baselines (B1 static, B2 static$+$mutation, B3 random-budget) and report them in
Table~\ref{tab:e1}. We also state explicitly in \S9 which claims remain unsupported. See
\chg{\S6.1 (Table~\ref{tab:e1}), \S9 (Table~\ref{tab:claim}).}

\rev{Why is mutation score not treated as defect detection?}
\resp Because the two measure different things; we cite Just et al.\ (FSE 2014) and show the gap
empirically (internal mutation core 97.3\% vs.\ external corpus 35.0\%). See \chg{\S7 item 4.}

\section*{Reviewer 2}

\rev{Reproducibility is unclear; the numbers depend on an undeclared environment.}
\resp We now declare the WSL + \texttt{g++} + \texttt{setarch} dependency in the abstract, \S9, and
Appendix~\ref{app:repro}, and require a verdict of \texttt{UNVERIFIED} (not a low score) when a
dependency is missing. See \chg{Abstract; \S9; Appendix~\ref{app:repro}.}

\rev{How does this differ from benchmark-evolution work?}
\resp Benchmark evolution changes the \emph{items}; we change the \emph{measuring stick} and require
the change to be externally measurable. See \chg{Table~\ref{tab:positioning}, \S2.}

\section*{Reviewer 3}

\rev{Small sample sizes undermine the statistical claims.}
\resp We agree; we report Clopper--Pearson CIs everywhere, mark $n<20$ rows as exploratory, and
state that $\Delta{=}15$pp needs $n\approx102$--172. See \chg{\S6, Appendix~\ref{app:stats}.}

\rev{Minor: terminology (detector/probe) is inconsistent.}
\resp Fixed; we use ``detector'' throughout. See \chg{whole manuscript.}

\section*{Standard replies (reusable blocks)}

\paragraph{On ``sample size is too small''.}
We agree, and we now say so in the abstract and \S9 rather than only in a limitations footnote.
Our holdout is $n{=}16$ measurable and the corpus $n{=}32$; $\pm$10pp of CI half-width would need
$n\approx96$. We therefore report \emph{direction} (the $\Delta$ lower bound is $\ge31.5$pp; exact
McNemar $p\le0.002$) and explicitly refuse to report magnitude. Sample expansion to $n\ge100$ is
listed as the first item of future work, with a written sampling plan.

\paragraph{On ``the baseline is not strong enough''.}
Correct, and we label it as such. The static arm is a \textbf{caliber re-binning}, not a re-run of a
real static detector; the random arm is an \textbf{instrument-level proxy}, not the paper's
budget-matched B3. Both are marked and excluded from the conclusion. The true B1/B3 require a
\texttt{detect\_static}/\texttt{select\_assets} interface that currently lives in a split
repository and is \textbf{blocked} for us; we register this as a blocked item rather than
substituting a convenient approximation.

\paragraph{On ``why not evaluate on SWE-bench''.}
Because our object of study is not code repair but \emph{verification of knowledge claims}. SWE-bench
measures whether a model can patch an issue; we measure whether a claim is true under an explicit
boundary triple. Adapting SWE-bench would change the construct (see \S8.1). We do cite the
SWE-bench contamination line as \emph{motivation} for external-sample-first design.

\paragraph{On ``reproducibility''.}
Every headline number is recomputed from landed artifacts by a single command; the environment
dependency (WSL + g++ + \texttt{setarch}) is declared, and a missing dependency yields
\texttt{UNVERIFIED} rather than a misleading low score. A five-tool paper pipeline (number sync,
BibTeX audit, figure-data provenance, anonymization, quality gate) runs in CI. Code is Apache-2.0
with DCO sign-off; Croissant metadata ships with the supplement.

\bigskip
\noindent
We hope the revised version addresses all concerns.

\end{document}
